\section{Introduction}\label{sec:introduction}

We determine the mixing exponent of single-site heat-bath dynamics
at the critical temperature of the zero-field Gaussian
Sherrington--Kirkpatrick model. Starting from the worst configuration,
the dynamics mixes in time $n^{2/3+o(1)}$ when each site has rate one,
or after $n^{5/3+o(1)}$ attempted uniform-site updates. The same
per-site exponent governs the relaxation time and the inverse classical
logarithmic Sobolev constant.

At criticality, equilibrium fluctuations are much larger than for
independent spins. An upper mixing bound must hold from every initial
configuration, while the functional bounds must hold for every test
function. We obtain these uniform controls in two ways. The first relates
a function's heat-bath energy to the mass of its support. The second proves
a classical logarithmic Sobolev inequality. Each argument has its own
initialization and entropy estimates.

Let $J^\circ$ be a symmetric matrix with zero diagonal and independent
$J^\circ_{ij}\sim N(0,1/n)$ for $i<j$. On $\Omega_n=\{-1,1\}^n$, put
\begin{equation}\label{eq:gibbs-law}
 \mu_J(x)=Z_J^{-1}\exp\!\left(\frac12x^{\mathsf T}J^\circ x\right).
\end{equation}
This is Gaussian SK at inverse temperature one and zero external field.
Write $Q_i$ for conditional expectation under $\mu_J$ given $x_{-i}$,
and set
\begin{equation}\label{eq:clock-conventions}
 \mathcal L_J=\sum_{i=1}^n(Q_i-I),\qquad
 P_J=\frac1n\sum_{i=1}^nQ_i,\qquad
 \D_J(f)=\sum_{i=1}^n\mu_J\!\left[\Var_{\mu_J}(f\mid x_{-i})\right].
\end{equation}
Thus $\mathcal L_J=n(P_J-I)$, and a discrete attempt may leave the spin
unchanged. Let $t_{\mathrm{mix},J}(\varepsilon)$ and
$k_{\mathrm{mix},J}(\varepsilon)$ denote worst-start total-variation
mixing times for $e^{t\mathcal L_J}$ and $P_J^k$, respectively.
The default threshold is $1/4$.

Define the relaxation time and the classical logarithmic Sobolev
inverse constant by
\begin{align}
 R_J&=\sup_{\Var_{\mu_J}(f)>0}\frac{\Var_{\mu_J}(f)}{\D_J(f)},
 \label{eq:relaxation-definition}\\
 C_{\mathrm{LS},J}&=\sup_{\Ent_{\mu_J}(f^2)>0}
               \frac{\Ent_{\mu_J}(f^2)}{\D_J(f)},
 \label{eq:lsi-definition}
\end{align}
where $\Ent_\mu(g)=\mu(g\log g)-\mu(g)\log\mu(g)$ for $g\ge0$.
Both constants use the unscaled sum of heat-bath conditional variances.
Thus $R_J$ is the relaxation time for the rate-one-per-site generator;
for one attempted update the relaxation time is $nR_J$. The conditional
variance weights in $\D_J$ are retained in every comparison below.

\Needspace{11\baselineskip}
\begin{theorem}[Critical mixing and functional inequalities]
\label{thm:main}
For every fixed $\varepsilon\in(0,1/2)$, in probability over the Gaussian
disorder,
\begin{align*}
 t_{\mathrm{mix},J}(\varepsilon)&=n^{2/3+o(1)},&
 k_{\mathrm{mix},J}(\varepsilon)&=n^{5/3+o(1)},\\
 R_J&=n^{2/3+o(1)},& C_{\mathrm{LS},J}&=n^{2/3+o(1)}.
\end{align*}
In each expression, $X_n=n^{a+o(1)}$ means that for every fixed
$\delta>0$, $\Pprob\{n^{a-\delta}\le X_n\le n^{a+\delta}\}\to1$.
\end{theorem}

The two functional constants describe complementary forms of relaxation.
For every centered observable $f$, the spectral theorem gives
$\Var_{\mu_J}(e^{t\mathcal L_J}f)\le e^{-2t/R_J}\Var_{\mu_J}(f)$,
with equality for a slowest eigenfunction. The classical logarithmic
Sobolev inequality supplies hypercontractive smoothing for arbitrary
densities and implies the Poincar\'e inequality. \Cref{sec:dynamics}
derives the worst-start bounds from these controls. Let $\operatorname{KL}$
denote relative entropy. \Cref{cor:entropy-production} also gives the
same per-site exponent for the inverse of the best rate $c$ such that
$\operatorname{KL}(\nu e^{t\mathcal L_J}\Vert\mu_J)
\le e^{-ct}\operatorname{KL}(\nu\Vert\mu_J)$
for every initial probability law $\nu$ on the cube and every $t\ge0$.

The contribution of this paper is the upper bounds. To state precisely
how they combine with the existing lower bounds, we record the critical
companion's realized-start and linear-observable theorems
\cite[Theorems~1.2 and~1.3]{AcceptedCritical}. Their quantifiers are
stronger than fixed exponent bounds and will be retained throughout.

\begin{theorem}[Realized equilibrium starts and linear observables]
\label{thm:accepted-critical}
For every deterministic sequence $t_n\ge0$ with $t_n=o(n^{2/3})$,
\[
 \mu_J\!\left\{x:
 \norm{e^{t_n\mathcal L_J}(x,\cdot)-\mu_J}_{\TV}>\frac14\right\}
 \longrightarrow1
\]
in probability over disorder. For every deterministic nonnegative integer
sequence $k_n=o(n^{5/3})$, the same assertion holds with $P_J^{k_n}$.
Let $\Sigma_J=\Cov_{\mu_J}(X)$ and
\[
 R_{\mathrm{lin},J}=\sup_{a\in\R^n\setminus\{0\}}
       \frac{\Var_{\mu_J}(a^{\mathsf T}X)}{\D_J(a^{\mathsf T}X)}.
\]
For every positive deterministic $M_n\to\infty$, with probability
tending to one both $\norm{\Sigma_J}$ and $R_{\mathrm{lin},J}$ belong to
$[n^{2/3}/M_n,M_n n^{2/3}]$.
\end{theorem}

In the realized-start assertion, one samples an initial configuration
from $\mu_J$ and holds it fixed before taking the transition law and its
TV distance. Averaging that transition law over the initial configuration
would instead give $\mu_J$ at every time. The last assertion concerns
covariance and linear Rayleigh quotients. Its arbitrary-diverging-factor
upper bound does not extend here to the full relaxation or logarithmic
Sobolev constant: the upper estimates in \Cref{thm:main} have every
fixed positive exponent slack.
For other fixed TV thresholds below $1/2$, the lower bound in
\Cref{thm:main} follows by testing the slowest gap eigenfunction, as
shown in \Cref{sec:dynamics}. The realized-start assertion itself
retains its stated threshold $1/4$.

The result concerns exactly the critical, zero-field Gaussian model.
It gives the mixing exponent and does not assert a cutoff, a cutoff
window, or a leading constant. The two upper-bound arguments retain
their separate intermediate conclusions and parameter choices.

\paragraph{The development of the dynamical question.}
Sherrington and Kirkpatrick introduced the infinite-range spin-glass
model in 1975 \cite{SherringtonKirkpatrick1975}. Their later analysis
with Gaussian couplings studied single-spin Glauber relaxation through
a linearized mean-field calculation and Monte Carlo simulations
\cite[Section~VII]{KirkpatrickSherrington1978}. A parallel dynamical
mean-field program described correlation and response functions in
soft-spin Langevin models; the infinite-range specialization was
developed by Sompolinsky and Zippelius
\cite{SompolinskyZippelius1982}. These studies examined relaxation near
the transition through correlation and response. The finite-size mixing
question asks how long the full distribution, started from the most
difficult configuration, takes to approach equilibrium.

Numerical work later examined how the critical time scale grows with the
number of spins. Billoire's study of slow samples used binary couplings
$J_{ij}=\pm n^{-1/2}$ \cite{Billoire2011SlowSamples}. Billoire and
Campbell's subsequent critical study measured the equilibrium
autocorrelation
$q(t)=n^{-1}\sum_i\langle X_i(0)X_i(t)\rangle$ after equilibration.
For heat-bath updates, with time measured in updates per spin, it tested
the finite-size scaling form
$q_c(n,t)\sim n^{-1/3}F(t/n^{2/3})$
\cite[equation~(7)]{BilloireCampbell2011CriticalDynamics}.
That critical preprint does not restate the coupling law. We use its
reported collapse as numerical evidence about equilibrium autocorrelation,
keeping it separate from the Gaussian worst-start question here.

\paragraph{Rigorous high-temperature dynamics.}
For comparison with the critical law \eqref{eq:gibbs-law}, inverse
temperature $\beta$ multiplies its exponent by $\beta$; the resulting
interaction entries have variance $\beta^2/n$. The SK statements below
hold with high probability over disorder. The inverse temperature
$\beta>0$ is fixed inside each stated range, and
mixing bounds use fixed total-variation accuracy unless $\varepsilon$
is displayed.
Spectral criteria gave functional inequalities well inside the
high-temperature phase. Bauerschmidt and Bodineau proved a logarithmic
Sobolev inequality under a spectral-width condition that applies to SK
for fixed $\beta<1/4$ \cite[Theorem~1 and Corollary~2]{BauerschmidtBodineau2019}.
For Ising spins their energy is the unweighted full spin-flip gradient.
The heat-bath form $\D_J$ instead retains conditional variance weights,
so that result is not the heat-bath inequality used in this paper.
Eldan, Koehler, and Zeitouni obtained a heat-bath spectral gap uniform
over external fields under the corresponding shifted spectral condition.
Their field-dependent worst-start bound gives polynomial mixing for
zero-field SK in the range $\beta<1/4$
\cite[Theorems~1 and~11, Section~5]{EldanKoehlerZeitouni2022}.
Anari, Jain, Koehler, Pham, and Vuong then obtained
$O_\beta(n\log n)$ random-site attempts in that range, uniformly over
external fields, through a modified logarithmic Sobolev inequality
\cite[Theorem~12]{AnariJainKoehlerPhamVuong2021}.

Subsequent work increased the temperature range. Anari, Koehler, and
Vuong derived a criterion governed by a
Volterra equation; it gives $O(n\log n)$ attempted updates for SK below
a reported numerical threshold of approximately $\beta=0.295$
\cite[Theorem~3 and Section~1.3]{AnariKoehlerVuong2024}.
Wang proved $O_\beta(n\log(n/\varepsilon))$ attempted updates for every
fixed $\beta<1/2$, with high probability simultaneously over external
fields and $\varepsilon\in(0,1)$
\cite[Theorem~1.1]{Wang2026OptimalMixing}.
At zero field, Boban, Li, and Oveis Gharan obtained a constant per-site
spectral gap and an $O(n^2)$ attempted-update bound for
$\beta<1/2+\varepsilon_0$, with a universal $\varepsilon_0>0$
\cite[Theorem~1.5 and Corollary~1.6]{BobanLiOveisGharan2026}.
The companion spectral-gap theorem covers each fixed zero-field
$0<\beta<1$, with a constant that may depend on $\beta$
\cite[Theorem~1.1]{AcceptedGap}. None of these fixed-subcritical bounds
supplies a bound at $\beta=1$ by taking a limit. At that endpoint the
critical companion identifies the slow linear and realized-start scales
in \Cref{thm:accepted-critical}; the present upper bounds match their
exponents for the full dynamics.

\paragraph{Static comparisons and the two critical mechanisms.}
The equilibrium geometry behind the critical scale is accessible through
comparison with spins on a sphere. Comets conditioned on the eigenvalues
of a GOE completion and averaged its eigenframe, obtaining an exact
conditional first-moment identity between cube and spherical partition
functions \cite[equations~(2.1)--(2.2)]{Comets1996Spherical}.
At criticality, Du and Huang developed sharper comparisons for free
energy and overlap. Their local overlap-cell estimates compare
partition-function moments
\cite[Theorem~1.6 and Section~3.3]{DuHuang2026CriticalFluctuations},
and their mixed cube--sphere tensor calculation cancels leading terms
in a squared norm
\cite[Section~3.2, Lemma~3.5]{DuHuang2026CriticalOverlap}.
These are static comparisons. The upper bounds here need corresponding
control after observations have changed the field and interaction,
together with probability estimates strong enough for arbitrary test
functions.

Gaussian stochastic localization supplies the framework for revealing
those observations while tracking the remaining variance or entropy
\cite{Eldan2013StochasticLocalization,ChenEldan2025Localization}.
The root calculation also develops conditional Kac--Rice methods used
for TAP landscapes by Fan, Mei, and Montanari, and by Celentano, Fan, and
Mei \cite{FanMeiMontanari2021,CelentanoFanMei2023}. Their cited
applications have different positive-overlap or fixed-signal hypotheses;
the shrinking-overlap critical estimates are proved here.
The isotropic route extends the signed spin--Gaussian comparison to
observation-dependent fields, then improves a typical-orientation estimate
along the observation path. The cap route keeps the reduction of covariance
in the field direction while fresh rotations reduce errors in orthogonal
directions inside the active spectral subspace. The entropy arguments
use these estimates only on their stated path or field domains and pay
for their exceptions. \Cref{sec:proof-overview} describes the inputs
and outputs of the two mechanisms in detail.
